% terminal-bench-cn — short paper draft (NeurIPS 2026 Datasets & Benchmarks track)
% Compile with `pdflatex main.tex` from this directory; `make pdf` from repo root also works.
\documentclass{article}

\usepackage[utf8]{inputenc}
\usepackage[T1]{fontenc}
\usepackage{xeCJK}        % Chinese support
\usepackage{microtype}
\usepackage{booktabs}
\usepackage{hyperref}
\usepackage{url}
\usepackage{graphicx}
\usepackage[margin=1in]{geometry}

\title{Terminal-Bench-CN: A Chinese-Language Slice for Evaluating LLM Coding Agents}
\author{ychenfen \\ \texttt{2570601904@qq.com}}
\date{\today}

\begin{document}
\maketitle

\begin{abstract}
We introduce Terminal-Bench-CN, a Chinese-language extension of Terminal-Bench
that evaluates coding agents on hermetic, dockerised tasks where the natural-language
ticket is delivered in Chinese. Each task ships a parallel English mirror, enabling
the first systematic study of cross-lingual deltas in coding-agent performance.
On v0.1 (5 tasks), we observe that frontier agents achieve XX\% pass-rate on
English tickets but YY\% on the Chinese ones — a ZZ-point gap concentrated in
encoding, path-handling and instruction-following failure modes. We release the
benchmark, the evaluation harness, and a cross-lingual delta analyzer at
\url{https://github.com/ychenfen/terminal-bench-cn}.
\end{abstract}

\section{Introduction}
% Motivation: Chinese is the primary work language for ~25\% of professional
% developers; existing coding-agent benchmarks (SWE-bench, Terminal-Bench,
% AgentBench) are English-only; translating prompts surfaces non-trivial failures.
\input{figs/leaderboard}

\section{Related Work}
% Coding-agent benchmarks: SWE-bench (Jimenez et al.), Terminal-Bench
% (Stanford NLP, 2025), AgentBench (Liu et al., 2023), AAMAS-2026 auto-benchmark
% line of work; cross-lingual NLP eval: XGLUE, XTREME-R; multilingual code:
% MBXP, Multilingual HumanEval (Athiwaratkun et al., 2023).

\section{Tasks}
The v0.1 release contains five hand-curated tasks spanning DevOps, encoding,
JavaScript toolchain, Django web back-end and TypeScript monorepos. Each task is
defined by:
\begin{itemize}
  \item \texttt{instruction.zh.md} — the Chinese ticket the agent reads
  \item \texttt{instruction.en.md} — a manually translated English mirror
  \item \texttt{docker/Dockerfile} — hermetic environment with the planted bug
  \item \texttt{tests/test\_NNN.sh} — a 0/1 grader; both correctness and "no-cheating" guards
\end{itemize}

\section{Method}
We run each agent under each language with $N=5$ samples per cell. Task scores
are 0/1 hermetic test pass; agent scores are micro-averaged across tasks.
We report cross-lingual delta $\Delta=p_{\rm zh}-p_{\rm en}$ with a 95\% Wilson interval
and bucket failure logs into a coarse 8-class taxonomy
(encoding, path/filename, dependency install, test assertion, config format,
agent gave up, ran out of steps, nuclear bypass).

\input{figs/delta_table}

\section{Results}
% Filled in by `make figures` from logs/cross_lingual_delta.json
% Discussion of patterns: encoding tasks dominate the zh-only failures;
% nuclear-bypass attempts higher on Chinese?

\section{Limitations and Future Work}
v0.1 is small. v0.2 will scale to 25 tasks and add traditional Chinese (zh-Hant);
v0.3 will host a public leaderboard and accept community contributions.

\section*{Acknowledgments}
Built on top of the excellent Terminal-Bench harness.

\bibliographystyle{plainnat}
\bibliography{refs}
\end{document}
